%----------------------------------------------------------------------------
\chapter*{Hallgatói nyilatkozat generatív mesterséges intelligencia alkalmazásáról}\label{aidecl}
\addcontentsline{toc}{chapter}{Hallgatói nyilatkozat generatív mesterséges intelligencia alkalmazásáról}
%----------------------------------------------------------------------------
% The faculty's form (docs/VIK_Generativ_MI_Hallgatoi_Nyilatkozat.docx), from
% https://vik.bme.hu/page/1346/ ; a filled-in example is in docs/VIK_Generativ_MI_Pelda_Nyilatkozat_Szakdolgozathoz.pdf.
% Keep the faculty's Hungarian wording; it is the form the faculty asks for.
% Included by draft.tex after the bibliography, before the appendices.
\todo[inline, caption={Fill in and sign the AI declaration.}]{Mark one box and complete every applicable row, the overall share and its justification. Record the tool and version, affected parts, estimated share, essential prompt, date and conversation link. Use -- for unused rows. Follow \guideref{Disclose AI use}, including the worked example and the citation instructions. For a thesis, sign the printed declaration.}

\selecthungarian
\newcommand{\aibox}{$\square$}        % unmarked box
\newcommand{\aiboxmarked}{$\boxtimes$} % marked box: swap it in for \aibox

\noindent\aibox\enspace Nem használtam semmilyen generatív MI segédeszközt.

\noindent\aibox\enspace Használtam generatív MI segédeszközt. Az MI-vel generált tartalmakat ellenőriztem, a generált kimenetek valóságtartalmáról meggyőződtem, az alábbi táblázatban megfelelően jelöltem minden használatot.

\newcommand{\aiuse}[1]{#1 & -- & -- & -- \\* \textit{Prompt lényegi része} & \multicolumn{3}{>{\raggedright\arraybackslash}p{\dimexpr0.66\linewidth+4\tabcolsep\relax}}{--} \\ \midrule}
{\small
\begin{longtable}{>{\raggedright\arraybackslash}p{0.24\linewidth} >{\raggedright\arraybackslash}p{0.17\linewidth} >{\raggedright\arraybackslash}p{0.31\linewidth} >{\raggedright\arraybackslash}p{0.14\linewidth}}
  \toprule
  \textbf{Felhasználási módok} & \textbf{Generatív MI eszköz(ök) neve} & \textbf{Érintett részek (fejezet, oldalszám, hivatkozás)} & \textbf{Használat becsült aránya (felhasználási módonként)} \\
  \midrule
  \endhead
  \aiuse{Irodalomkutatás}
  \aiuse{Programkód generálása}
  \aiuse{Új ötletek, megoldási javaslatok generálása}
  \aiuse{Vázlat létrehozása (szövegstruktúra, vázlatpontok)}
  \aiuse{Szövegblokkok létrehozása}
  \aiuse{Képek generálása illusztrációs célból}
  \aiuse{Adatvizualizáció, grafikonok generálása adatpontok alapján}
  \aiuse{Prezentáció készítése}
  \aiuse{Egyéb (nevezze meg)}
  \multicolumn{3}{l}{\textbf{Összesített százalékos érték (a feladat érdemi részére nézve)}} & -- \\
  \midrule
  \multicolumn{4}{>{\raggedright\arraybackslash}p{\dimexpr0.86\linewidth+6\tabcolsep\relax}}{\textbf{Összesített érték rövid, szöveges indoklása:} TODO} \\
  \bottomrule
\end{longtable}
}

\noindent\textbf{Figyelem:} a fenti százalékértékektől függetlenül az MI-vel generált minden tartalomért a szerző a felelős!

\bigskip
% A thesis is signed on paper; a TDK paper is uploaded as a PDF, so it gets the date and the name without a signature line.
\expandafter\ifx\vikdoktipus\tdk
  \noindent Kelt: Budapest, \today\hfill\szerzo
\else
  \noindent Kelt: Budapest, \today\hfill\begin{tabular}[t]{c}\rule{5cm}{0.4pt}\\\szerzo\end{tabular}
\fi

\selectthesislanguage
